For exercises 3.1 and 3.5, determine the range (possible values) of the random variable.

\enquestion{3.1}{
    The random variable is the number of nonconforming solder connections on a printed circuit board with $2000$ connections.
}

\enanswer{
    The range of this discrete random variable is the set $\set{0, 1, 2, \dots, 2000}$, as each of the $2000$ connections can be nonconforming or not.
}

\enquestion{3.5}{
    A batch of $600$ machined parts contains $10$ that do not conform to customer requirements. 
    Parts are selected successively, without replacement, until a nonconforming part is obtained.
    The random variable is the number of parts selected.
}

\enanswer{
    Notice that we keep picking until a nonconforming part is obtained, so the lower bound is $1$ if we immediately pick one,
    and can take up to $591$ tries if we have bad luck and first take all $590$ conforming parts before we pick a nonconforming one.
    Therefore, the range of this discrete random variable is the set $\set{1, 2, 3, \dots, 591}$.
}

\enquestion{3.14}{
    The sample space of a random experiment is $\set{a,b,c,d,e,f}$, and each outcome is equally likely.
    A random variable is defined as follows:
    \begin{center}
        \begin{tabular}{*{7}{c}}
            \toprule
                \textbf{Outcome}& $a$ & $b$ & $c$ & $d$ & $e$ & $f$ \\
            \midrule
                $x$ & $0$ & $0$ & $1.5$ & $1.5$ & $2$ & $3$ \\
            \bottomrule
        \end{tabular}
    \end{center}
    Determine the probability mass function of $X$.
    Use the probability mass function to determine the following probabilities:

    \begin{tabular}{*{2}{m{.45\textwidth}}}
        (a) $P(X=2)$ & (b) $P(0.6 < X < 2.7)$ \\
        (c) $P(X>3)$ & (d) $P(0 < X \leq 2)$ \\
        (e) $P(X=0 \text{ or } X=2)$ & \\
    \end{tabular}
}

\enanswer{
    Each of the outcomes of the random experiment occurs with equal probability, so we can use the classical approach to probability theory:
    \begin{enumerate}[label=(\alph*)]
        \item $P(X=2) = \frac{\text{number of favorable outcomes}}{\text{total number of outcomes}} = \frac{1}{6} \approx 0.1667$,
        \item $P(0.6 < X < 2.7) = \frac{\text{number of favorable outcomes}}{\text{total number of outcomes}} = \frac{3}{6} = 0.5$,
        \item $P(0.6 < X < 2.7) = \frac{\text{number of favorable outcomes}}{\text{total number of outcomes}} = \frac{3}{6} = 0.5$,
        \item $P(0 < X \leq 2) = \frac{\text{number of favorable outcomes}}{\text{total number of outcomes}} = \frac{5}{6} \approx 0.8333$,
        \item $P(X=0 \text{ or } X=2) = \frac{\text{number of favorable outcomes}}{\text{total number of outcomes}} = \frac{3}{6} = 0.5$.
    \end{enumerate}
}


\enquestion{3.17}{
    Verify that the following function is a probability mass function, and determine the requested probabilities.
    \[
        f(x) = \frac{2x+1}{25}, \qquad x=0,1,2,3,4
    \]
    \begin{tabular}{*{2}{m{.45\textwidth}}}
        (a) $P(X=3)$            & (b) $P(X \leq 1)$ \\
        (c) $P(2 \leq X < 4)$ & (d) $P(X > -10)$ \\
    \end{tabular}
}

\enanswer{
    This function is a probability mass function, since for each possible value of $x$ the function value $f(x) \geq 0$.
    Furthermore, the sum for $x=0,1,2,3,4$ of $f(x)$ equals $1$.
    \begin{enumerate}[label=(\alph*)]
        \item $P(X=3) = f(3) = \frac{2 \cdot 3 + 1}{25} = \frac{7}{25} = 0.28$.
        \item $P(X\leq 1) = f(0) + f(1) = \frac{2 \cdot 0 + 1}{25} + \frac{2 \cdot 1 + 1}{25} = \frac{1}{25} + \frac{3}{25} = 0.16$.
        \item $P(2 \leq X < 4) = f(2) + f(3) = \frac{2 \cdot 2 + 1}{25} + \frac{2 \cdot 3 + 1}{25} = \frac{5}{25} + \frac{7}{25} = 0.48$.
        \item All the probability mass lies on outcomes $x=0,1,2,3,4$, which are larger than $10$. Therefore: $P(X > -10) = f(0) + f(1) + f(2) + f(3) + f(4) = 1$.
    \end{enumerate}
}

\enquestion{3.29}{
    An assembly consists of three mechanical components. 
    Suppose that the probabilities that the first, second, and third components meet specifications are $0.95$, $0.98$, and $0.97$, respectively.
    Assume that the components are independent.
    Determine the probability mass function of the number of components in the assembly that meet specifications.
}

\enanswer{
    Define the events $M_1$, $M_2$ and $M_3$ indicating whether the three components meet specifications.
    Then $P(M_1) = 0.95$, $P(M_2) = 0.98$ and $P(M_3) = 0.97$.
    Furthermore, the events are independent.
    Then, the number of components meeting specifications has the following probability mass function:
    \begin{itemize}
        \item $k = 0$: neither of the three events occur, so $f(0) = P(M_1' \cap M_2' \cap M_3') = P(M_1') \cdot P(M_2') \cdot P(M_3') = 0.05 \cdot 0.02 \cdot 0.03  = 0.00003$.
        \item $k = 1$: exactly one of the three events occur:
        \begin{align*}
            f(1)    &= P(M_1 \cap M_2' \cap M_3') +  P(M_1' \cap M_2 \cap M_3') + P(M_1' \cap M_2' \cap M_3) \\
                    &= 0.95 \cdot 0.02 \cdot 0.03 + 0.05 \cdot 0.98 \cdot 0.03 + 0.05 \cdot 0.02 \cdot 0.97 \\
                    &= 0.00301.
        \end{align*}
        \item $k = 2$: exactly two of the three events occur:
        \begin{align*}
            f(2)    &= P(M_1 \cap M_2 \cap M_3') +  P(M_1 \cap M_2' \cap M_3) + P(M_1' \cap M_2 \cap M_3) \\
                    &= 0.95 \cdot 0.98 \cdot 0.03 + 0.95 \cdot 0.02 \cdot 0.97 + 0.05 \cdot 0.98 \cdot 0.97 \\
                    &= 0.09389.
        \end{align*}
        \item $k = 3$: all three events occur:
        \begin{align*}
            f(3)    &= P(M_1 \cap M_2 \cap M_3) = 0.95 \cdot 0.98 \cdot 0.97 = 0.90307.
        \end{align*}
    \end{itemize}
    Summarized in a table, we have:

    \begin{tabular}{*{5}{m{0.15\textwidth}}}
        \toprule
            $x$ & $0$ & $1$ & $2$ & $3$ \\
        \midrule
            $f(x)$ & $0.00003$ & $0.00301$ & $0.09389$ & $0.90307$ \\
        \bottomrule
    \end{tabular}
}

\enquestion{3.35}{
    Determine the cumulative distribution function of the random variable with the following probability mass function:
    \begin{center}
        \begin{tabular}{*{6}{c}}
            \toprule
                $x$ & $-2$ & $-1$ & $0$ & $1$ & $2$ \\
            \midrule
                $f(x)$ & $1/8$ & $2/8$ & $2/8$ & $2/8$ & $1/8$ \\
            \bottomrule        
        \end{tabular}
    \end{center}
    Furthermore, determine the following probabilities:

    \begin{tabular}{*{2}{m{.45\textwidth}}}
        (a) $P(X \leq 1.25)$ & (b) $P(X \leq 2.2)$ \\
        (c) $P(-1.1 < X \leq 1)$ & (d) $P(X > 0)$ \\
    \end{tabular}
}

\enanswer{
    We can determine the cumulative distribution function (CDF) $F(x)$ by adding probabilities from left to right:
    \begin{center}
        \begin{tabular}{*{6}{c}}
            \toprule
                $x$ & $-2$ & $-1$ & $0$ & $1$ & $2$ \\
            \midrule
                $f(x)$ & $1/8$ & $2/8$ & $2/8$ & $2/8$ & $1/8$ \\
            \midrule
                $F(x)$ & $1/8$ & $3/8$ & $5/8$ & $7/8$ & $1$ \\
            \bottomrule        
        \end{tabular}
    \end{center}
    From this table, we can read the following probabilities:
    \begin{enumerate}
        \item $P(X \leq 1.25) = P(X \leq 1) = \frac{7}{8} = 0.875$, 
        \item $P(X \leq 2.2) = P(X \leq 2) = 1$,
        \item $P(-1.1 < X \leq 1) = P(X \leq 1) - P(X \leq -1.1) = \frac{7}{8} - \frac{1}{8} = 0.75$,
        \item $P(X > 0) = 1 - P(X \leq 0) = 1 - \frac{5}{8} = 0.375$.
    \end{enumerate}
}

\enquestion{3.43}{
    Verify that the following function is a cumulative distribution function, and determine the probability mass function and the requested probabilities:
    \[
        F(x) = \begin{cases}
            0 & x < -10 \\
            0.3 & -10 \leq x < 30 \\
            0.7 & 30 \leq x < 50 \\
            1 & 50 \leq x
        \end{cases}
    \]
    \begin{tabular}{*{2}{m{.45\textwidth}}}
        (a) $P(X \leq 50)$ & (b) $P(X \leq 40)$ \\
        (c) $P(40 \leq X \leq 60)$ & (d) $P(X < 0)$ \\
        (e) $P(0 \leq X < 10)$ & (f) $P(-10 < X < 10)$ \\
    \end{tabular}
}

\enanswer{
    We see that the function $F(x)$ is $0$ for small $x$ and $1$ for large $x$, and it is monotonic non-decreasing (as going to the right we add more and more probabilities).
    The probability mass function is given by taking the boundary values for $x$ for distinct cases and take the difference in $F(x)$:
    \begin{center}
        \begin{tabular}{*{6}{c}}
            \toprule
                $x$ & $-10$ & $30$ & $50$ \\
            \midrule
                $f(x)$ & $0.3$ & $0.4$ & $0.3$ \\
            \bottomrule        
        \end{tabular}
    \end{center}

    From this, we can derive the probabilities
    \begin{enumerate}
        \item $P(X \leq 50) = F(50) = 1$,
        \item $P(X \leq 40) = P(X \leq 30) = F(30) = 0.7$,
        \item $P(40 \leq X \leq 60) = P(X = 50) = f(50) = 0.3$,
        \item $P(X < 0) = P(X = -10) = f(-10) = 0.3$,
        \item $P(0 \leq X < 10) = 0$,
        \item $P(-10 < X < 10) = 0$.
    \end{enumerate}
}


\enquestion{3.51}{
    Determine the mean and variance of the following random variable:
    \begin{center}
        \begin{tabular}{*{6}{c}}
            \toprule
                $x$ & $-2$ & $-1$ & $0$ & $1$ & $2$ \\
            \midrule
                $f(x)$ & $1/8$ & $2/8$ & $2/8$ & $2/8$ & $1/8$ \\
            \bottomrule        
        \end{tabular}
    \end{center}
}

\enanswer{
    We can determine the mean of the random variable $X$ as follows:
    \begin{align*}
        E(X)    &= \sum_{x=-2}^{2} x \cdot f(x) \\
                &= -2 \cdot 1/8 + -1 \cdot 2/8 + 0 \cdot 2/8 + 1 \cdot 2/8 + 2 \cdot 1/8 \\
                &= 0.
    \end{align*}
    We can determine the variance of the random variable $X$ as follows:
    \begin{align*}
        \Var{X} &= \sum_{x=-2}^{2} (x - E(X))^2 \cdot f(x) \\
                &= (-2-0)^2 \cdot 1/8 + (-1-0)^2 \cdot 2/8 + (0-0)^2 \cdot 2/8 + (1-0)^2 \cdot 2/8 + (2-0)^2 \cdot 1/8 \\
                &= 1.5.
    \end{align*}
}


\enquestion{3.53}{
    Determine the mean and variance of the following random variable with the following probability mass function:
    \[
        f(x) = \frac{2x+1}{25}, \qquad x=0,1,2,3,4
    \]
}

\enanswer{
    We can determine the mean of the random variable $X$ as follows:
    \begin{align*}
        E(X)    &= \sum_{x=0}^{4} x \cdot f(x) \\
                &= 0 \cdot 1/25 + 1 \cdot 3/25 + 2 \cdot 5/25 + 3 \cdot 7/25 + 4 \cdot 9/25 \\
                &= 2.8.
    \end{align*}
    We can determine the variance of the random variable $X$ as follows:
    \begin{align*}
        \Var{X} &= \sum_{x=0}^{4} (x - E(X))^2 \cdot f(x) \\
                &= (0-2.8)^2 \cdot 1/25 + (1-2.8)^2 \cdot 3/25 + (2-2.8)^2 \cdot 5/25 + (3-2.8)^2 \cdot 7/25 + (4-2.8)^2 \cdot 9/25 \\
                &= 1.36.
    \end{align*}
}

\enquestion{3.67}{
    Let the random variable $X$ have a discrete uniform distribution on the integers $1 \leq x \leq 5$.
    Determine the mean and variance of $X$.
}

\enanswer{
    The discrete uniform distribution on the integers $1 \leq x \leq 5$ has the following probability mass function:
    \begin{center}
        \begin{tabular}{*{6}{c}}
            \toprule
                $x$ & $1$ & $2$ & $3$ & $4$ & $5$ \\
            \midrule
                $f(x)$ & $1/5$ & $1/5$ & $1/5$ & $1/5$ & $1/5$ \\
            \bottomrule        
        \end{tabular}
    \end{center}

    We can determine the mean of the random variable $X$ as follows:
    \begin{align*}
        E(X)    &= \sum_{x=1}^{5} x \cdot f(x) \\
                &= 1 \cdot 1/5 + 2 \cdot 1/5 + 3 \cdot 1/5 + 4 \cdot 1/5 + 5 \cdot 1/5 \\
                &= 3.
    \end{align*}
    We can determine the variance of the random variable $X$ as follows:
    \begin{align*}
        \Var{X} &= \sum_{x=1}^{5} (x - E(X))^2 \cdot f(x) \\
                &= (1-3)^2 \cdot 1/5 + (2-3)^2 \cdot 1/5 + (3-3)^2 \cdot 1/5 + (4-3)^2 \cdot 1/5 + (5-3)^2 \cdot 1/5 \\
                &= 2.
    \end{align*}
}

\enquestion{3.68}{
    Thickness measurements of a coating process are made to the nearest hundredth of a millimeter.
    The thickness measurements are uniformly distributed with values $0.14$, $0.15$, $0.16$, $0.17$, $0.18$, $0.19$, and $0.20$.
    Determine the mean and variance of the coating thickness for this process.
}

\enanswer{
    The discrete uniform distribution on these values has the following probability mass function:
    \begin{center}
        \begin{tabular}{*{8}{c}}
            \toprule
                $x$ & $0.14$ & $0.15$ & $0.16$ & $0.17$ & $0.18$ & $0.19$ & $0.20$ \\
            \midrule
                $f(x)$ & $1/7$ & $1/7$ & $1/7$ & $1/7$ & $1/7$ & $1/7$ & $1/7$ \\
            \bottomrule        
        \end{tabular}
    \end{center}

    We can determine the mean of the random variable $X$ as follows:
    \begin{align*}
        E(X)    &= \sum_{x} x \cdot f(x) \\
                &= 0.14 \cdot 1/7 + 0.15 \cdot 1/7 + \ldots + 0.20 \cdot 1/7 \\
                &= 0.17.
    \end{align*}
    We can determine the variance of the random variable $X$ as follows:
    \begin{align*}
        \Var{X} &= \sum_{x} (x - E(X))^2 \cdot f(x) \\
                &= (0.14-0.17)^2 \cdot 1/7 + (0.15-0.17)^2 \cdot 1/7 + \ldots + (0.20-0.17)^2 \cdot 1/7 \\
                &= 0.0004.
    \end{align*}
}

\enquestion{3.70}{
    The lengths of plate glass parts are measured to the nearest tenth of a millimeter.
    The lengths are uniformly distributed with values at every tenth of a millimeter starting at $640.0$ and continuing through $640.9$.
    Determine the mean and variance of the lengths.
}

\enanswer{
    The discrete uniform distribution on these values has the following probability mass function:
    \begin{center}
        \begin{tabular}{*{11}{c}}
            \toprule
                $x$ & $640.0$ & $640.1$ & $640.2$ & $640.3$ & $640.4$ & $640.5$ & $640.6$ & $640.7$ & $640.8$ & $640.9$ \\
            \midrule
                $f(x)$ & $1/10$ & $1/10$ & $1/10$ & $1/10$ & $1/10$ & $1/10$ & $1/10$ & $1/10$ & $1/10$ & $1/10$ \\
            \bottomrule        
        \end{tabular}
    \end{center}

    We can determine the mean of the random variable $X$ as follows:
    \begin{align*}
        E(X)    &= \sum_{x} x \cdot f(x) \\
                &= 640.0 \cdot 1/10 + 640.1 \cdot 1/10 + \ldots + 640.9 \cdot 1/10 \\
                &= 640.45.
    \end{align*}
    We can determine the variance of the random variable $X$ as follows:
    \begin{align*}
        \Var{X} &= \sum_{x} (x - E(X))^2 \cdot f(x) \\
                &= (640.0 - 640.45)^2 \cdot 1/10 + (640.1 - 640.45)^2 \cdot 1/10 + \ldots + (640.9-640.45)^2 \cdot 1/10 \\
                &= 0.0825.
    \end{align*}
}